\documentclass{zhvt-classic}

% 使用本地源流明体（GenRyuMin2 TW）按文件路径加载，不依赖系统字体名
% 版框（内框／界栏／书口侧栏／鱼尾）改用与版心印章同一印泥棕色：
%   印章图片 yshs.png 实色 = RGB(135,68,52) = HTML 874434
\usepackage{xcolor}
\usepackage{tikz}
% sealbrown 双注册：xcolor 供 TikZ/pgf，L3 供 cls 版框
\definecolor{sealbrown}{HTML}{874434}
\ExplSyntaxOn
\color_set:nnn { sealbrown } { HTML } { 874434 }
\ExplSyntaxOff

\zhvtset{
  grid color       = sealbrown,
  main font        = GenRyuMin2TW-SB,
  main font path   = C:/Users/Administrator/texmf/fonts/opentype/public/GenRyuMin2TW/,
jiazhu font      = GenRyuMin2TW-M,
  jiazhu font path = C:/Users/Administrator/texmf/fonts/opentype/public/GenRyuMin2TW/,
  % 页码框容纳 4 字：正常页「九十七」+前置空 1 格=4 行恰好；
  % 100+ 页「一百零一」共 5 行超出，末字「一/二/三」被框裁掉（2026-10-04 实测），故放宽到 5
  shukou page number chars = 5,
  % 书口书名框高从默认 6 字减到 4 字，把省出的 2 字高拨给 chapter 框，
  % 使 6 字书名+卷名放进 9 字高的 chapter 框后不再压线（2026-10-05）
  shukou title chars = 4,
}

\title{史记}[韩世家]
\author{汉·司马迁}
\maker{闲叟漫记·重製}


% ===== 索隐述赞样式（2026-11-08 四次定稿 · 愚大人亲手改定参数）=====
% ⚠ 竖排前提：自定义字体家族必须带 RawFeature={vertical:+vert}；
%   simhei.ttf 加 +vhal 会整体加载失败静默回退，只写 +vert；simfang 同理。
\newCJKfontfamily{\shuzanfont}[Path=C:/Windows/Fonts/,RawFeature={vertical:+vert}]{simhei.ttf}
\newCJKfontfamily{\shuzanbody}[Path=C:/Windows/Fonts/,RawFeature={vertical:+vert}]{simfang.ttf}
% 题头：152x30pt 圆角 6pt 棕底白字；baseline (0,2) 使整框落入界行格子不压版框线；
% 白字节点 x 不补偿、y +8.2pt（像素级墨迹实测 dx/dy≈0，L/R 留白约 20pt）。
\newcommand{\shuzanhead}[1]{%
  \begin{tikzpicture}[baseline={(0,2)}]
    \fill[sealbrown, rounded corners=6pt]
      (-76pt,-15pt) rectangle (76pt,15pt);
    \node[text=white, inner sep=0pt, anchor=center] at (0pt,8.2pt)
      {\shuzanfont\fontsize{23}{23}\selectfont #1};
  \end{tikzpicture}%
}
% ★ 修正书口卷名位置：chapter 框起始位置改回 cls 默认值，避免全局影响。
%   只在每个卷首页通过 \juan 命令里的空命令微调该页位置（安全，不会跨页污染）。
\ExplSyntaxOn
\AtBeginDocument{
  \dim_set:Nn \l__zhvt_shukou_chapter_title_start_dim
    { \l__zhvt_shukou_title_height_dim + \l__zhvt_font_size_dim * 2 }
  \dim_set:Nn \l__zhvt_shukou_page_number_start_dim
    { \l__zhvt_shukou_chapter_title_start_dim + \l__zhvt_shukou_chapter_title_height_dim }
}
\ExplSyntaxOff

\begin{document}

% ★★ 版心（书口）排布覆写：把书名放到鱼尾之下（2026-10-04）
%
%   cls 默认版心自上而下：书名（title 静态框，在鱼尾之上）→ 鱼尾 → 篇名（chapter 动态框）→ 页码
%   本模板改为：          鱼尾（位置不动）→ 书名 → 篇名 → 页码分割线 → 页码
%
%   手法：把 cls 放在 title 静态框里的书名，改放进 chapter 动态框（鱼尾之下）；
%         title 静态框留空。书名取 \l__zhvt_title_str，即 \title{...} 的第一参数，
%         所以换书时只改 \title，不必在此重复写书名。
%
%   ⚠ 三条必须遵守的写法（2026-10-04 实测）：
%     ① 汉字必须用 \hbox{...} 包住。flowfram 会把框内容装进一个宽 14.4mm（= 一个 line_sep
%        = 书口栏宽）的纵向 minipage 并按该宽度折行，而整页靠 PDF /Rotate 90 显示、
%        框内文字本身不变向 —— 不包 \hbox 时「围炉夜话」4 字会被逐字折成 4 行，
%        在版面上就成了横排 4 字并横向溢出版框。
%     ② 竖排下**一个 \hbox 就是一整列**。书名与篇名要同处一列，必须放进同一个 \hbox；
%        分成两个 \hbox 时，后一个会被当成新的一列、甩到隔壁栏去
%        （实测「序」跑到书口栏外侧 14.4mm 的正文列上）。中间的 \quad 在框内是横向空白，
%        显示出来正好是书名与篇名之间的纵向间隔。
%     ③ ★ 内容盒要写成 **与框等宽的定宽盒** `\hbox to \l__zhvt_shukou_chapter_title_height_dim{…\hfil}`。
%        实测：chapter 框是把内容盒沿纵向**居中**摆放的，内容宽度一变，书名就跟着漂：
%        自然宽度的 \hbox（书名 38.4mm）落在 x=157.2mm，加篇名后（57.6mm）落在 x=152.4mm，
%        目录页（无篇名）与正文页的书名就差半格。做成等宽盒后，无论有没有篇名，
%        内容固定从框顶起排，实测两页均落在 x=142.8mm，不再漂。
\ExplSyntaxOn
% 把书口 chapter 框（书名+篇名）起始位置下移 1 字高，避免 4 字卷名顶到鱼尾
% ⚠ cls 的 \__zhvt_calculate_dimensions: 在 \AtBeginDocument 里执行，会重算该 dim；
%   必须把覆写挂到 begin document 之后，否则会被二次覆盖（2026-10-05 实测）。
% (旧方案已移到导言区，直接修改 chapter 框起始位置)
% 卷标题命令 #1=书口篇名(鱼尾) #2=正文卷名（含卷序，卷名内空格用 ~）
%
% ★★ 目录页两项修订（2026-10-05，愚大人指出）★★
%   ① 页码起点统一对齐：
%      原状 —— 目录项写成 \par\hspace{4em}卷名\hspace{4em}页码，页码起点 = 卷名实宽 + 4em。
%      卷名实宽随字数变：「卷一~梁惠王上」6 字 +~、「卷七~离娄上」5 字 +~，
%      差 1 个字位 → 卷七~卷十 4 条的页码整体比其余 10 条高 27.21pt（一个字），看着上下不齐。
%      改法 —— 卷名补全角空格（宽 = 1 字 = 27.215pt）补足到 **7 字符位**再写进目录，
%      14 条页码起点遂全部落在同一条竖线上（起点对齐；末字随页码字数自然参差）。
%      ⚠ 只对目录里的副本 \l_tmpc_tl 补，正文卷名仍用 \l_tmpa_tl 原样 —— 正文像素不受影响。
%   ② 卷名整体上移半字位：前导 \hspace 由 4em 改 3.5em。
%      实测原状目录块在版框内偏下（天头留白 139.6pt / 地脚留白 110.7pt），
%      上移 0.5em = 13.6pt 后上下基本均分，视觉居中。
\NewDocumentCommand { \juan } { mm } {
	\clearpage
	\tl_set:Nn \l__zhvt_chapter_title_tl { #1 }
	\tl_set:Nn \l_tmpa_tl { #2 }
	\tl_set_eq:NN \l_tmpc_tl \l_tmpa_tl
	\int_step_inline:nn { 7 - \str_count:n { #2 } }
	  { \tl_put_right:Nn \l_tmpc_tl { 　 } }
	\tl_set:Nx \l_tmpb_tl { \zhnum { page } }
	\tl_put_left:Nn \l_tmpb_tl { \hspace { 4em } }
	\tl_put_left:Nn \l_tmpb_tl { \l_tmpc_tl }
	\tl_put_left:Nn \l_tmpb_tl { \par \hspace { 3.5em } }
	\addtocontents { toc } { \l_tmpb_tl }
	\skip_horizontal:n { \l__zhvt_font_size_dim * 4 }
	\l_tmpa_tl
	\par
	\thispagestyle { empty }
}
\cs_set_protected:Npn \__zhvt_set_shukou_content: {
	% ★ 鱼尾上方留白（title 静态框）嵌入闲叟钤印：
	%   整页 /Rotate 90，故此处把图用 \rotatebox{90} 转正为竖章；
	%   尺寸卡在版心条幅内（书内约 7mm 宽 × 33.8mm 高，留余量不压鱼尾）。
	\setstaticcontents * { title } {
		\vbox_to_ht:nn { \l__zhvt_baseline_skip_dim } {
			\vskip 14.2pt
			\hbox_to_wd:nn { \l__zhvt_shukou_title_height_dim } {
				\hfil
				\rotatebox [ origin = c ] { 90 } { \includegraphics [ width = 7mm , height = 33.8mm ] { yshs.png } }
				\hfil
			}
		}
	}
% ★ 版心（书口）统一改法（2026-11-16 愚大人定）：18pt＋框内居中＋书口文字同版框色
%   ① 原 27.2pt 时 7 字长篇名（如「廉颇蔺相如列传」）末字溢出 chapter 框、压住页码；
%      改 \fontsize{18pt}{18pt}，前后 \hfil 令字串在框内居中（不随篇名字数漂移）。
%   ② \vbox{\vskip 8.2pt} 定量下移，抵消缩字号后栏宽方向 8.2pt 上偏（像素实测 y 中心 543.3 = 原版）。
%      ⚠ \raisebox / \vspace* 在 flowfram 框内无效；换字号须同步换补偿：16pt→10.1、19pt→7.4。
%   ③ chapter 与 page~number 均置 \color{sealbrown}（HTML 874434），与版框／界行同色；正文仍黑。
%   ⚠ 页码字位、钤印、鱼尾、界行、正文一律未动；五部合订本同此口径。
	\setdynamiccontents * { chapter }
	{ \color{sealbrown} \vbox { \vskip 8.2pt \hbox to \l__zhvt_shukou_chapter_title_height_dim
		{ \hfil \fontsize{18pt}{18pt}\selectfont \l__zhvt_title_str \quad \l__zhvt_chapter_title_tl \hfil } } }
	\setdynamiccontents * { page~number } { \color{sealbrown} \quad \zhnum { page } }
}
\__zhvt_set_shukou_content:
\ExplSyntaxOff

\maketitle

% 单篇，无多卷目录，不渲染目录页

\juan{韩世家}{卷四十五~韩世家第十五}
\hspace*{2em}韩之先与周同姓。姓姬氏。其后苗裔事晋。得封于韩原。曰韩武子。武子后三世有韩厥。从封姓为韩氏。\\
\hspace*{2em}韩厥。晋景公之三年。晋司寇屠岸贾将作乱。诛灵公之贼赵盾。赵盾已死矣。欲诛其子赵朔。韩厥止贾。贾不听。厥告赵朔令亡。朔曰。﹁子必能不绝赵祀。死不恨矣。﹂韩厥许之。及贾诛赵氏。厥称疾不出。程婴。公孙杵臼之藏赵孤赵武也。厥知之。\\
\hspace*{2em}景公十一年。厥与郄克将兵八百乘伐齐。败齐顷公于鞍。获逢丑父。于是晋作六卿。而韩厥在一卿之位。号为献子。\\
\hspace*{2em}晋景公十七年。病。卜大业之不遂者为祟。韩厥称赵成季之功。今后无祀。以感景公。景公问曰。﹁尚有世乎。﹂厥于是言赵武。而复与故赵氏田邑。续赵氏祀。\\
\hspace*{2em}晋悼公之七年。韩献子老。献子卒。子宣子代。宣子徙居州。\\
\hspace*{2em}晋平公十四年。吴季札使晋。曰。﹁晋国之政卒归于韩。魏。赵矣。﹂晋顷公十二年。韩宣子与赵。魏共分祁氏。羊舌氏十县。晋定公十五年。宣子与赵简子侵伐范。中行氏。宣子卒。子贞子代立。贞子徙居平阳。\\
\hspace*{2em}贞子卒。子简子代。简子卒。子庄子代。庄子卒。子康子代。康子与赵襄子。魏桓子共败知伯。分其地。地益大。大于诸侯。\\
\hspace*{2em}康子卒。子武子代。武子二年。伐郑。杀其君幽公。十六年。武子卒。子景侯立。\\
\hspace*{2em}景侯虔元年。伐郑。取雍丘。二年。郑败我负黍。\\
\hspace*{2em}六年。与赵。魏俱得列为诸侯。\\
\hspace*{2em}九年。郑围我阳翟。景侯卒。子列侯取立。\\
\hspace*{2em}列侯三年。聂政杀韩相侠累。九年。秦伐我宜阳。取六邑。十三年。列侯卒。子文侯立。是岁魏文侯卒。\\
\hspace*{2em}文侯二年。伐郑。取阳城。伐宋。到彭城。执宋君。七年。伐齐。至桑丘。郑反晋。九年。伐齐。至灵丘。十年。文侯卒。子哀侯立。\\
\hspace*{2em}哀侯元年。与赵。魏分晋国。二年。灭郑。因徙都郑。\\
\hspace*{2em}六年。韩严弑其君哀侯。而子懿侯立。\\
\hspace*{2em}懿侯二年。魏败我马陵。五年。与魏惠王会宅阳。九年。魏败我浍。十二年。懿侯卒。子昭侯立。\\
\hspace*{2em}昭侯元年。秦败我西山。二年。宋取我黄池。魏取朱。六年。伐东周。取陵观。邢丘。\\
\hspace*{2em}八年。申不害相韩。修术行道。国内以治。诸侯不来侵伐。\\
\hspace*{2em}十年。韩姬弑其君悼公。十一年。昭侯如秦。二十二年。申不害死。二十四年。秦来拔我宜阳。\\
\hspace*{2em}二十五年。旱。作高门。屈宜臼曰。﹁昭侯不出此门。何也。不时。吾所谓时者。非时日也。人固有利不利时。昭侯尝利矣。不作高门。往年秦拔宜阳。今年旱。昭侯不以此时恤民之急。而顾益奢。此谓﹃时绌举赢﹄。﹂二十六年。高门成。昭侯卒。果不出此门。子宣惠王立。\\
\hspace*{2em}宣惠王五年。张仪相秦。八年。魏败我将韩举。十一年。君号为王。与赵会区鼠。十四年。秦伐败我鄢。\\
\hspace*{2em}十六年。秦败我修鱼。虏得韩将。申差于浊泽。韩氏急。公仲谓韩王曰。﹁与国非可恃也。今秦之欲伐楚久矣。王不如因张仪为和于秦。赂以一名都。具甲。与之南伐楚。此以一易二之计也。﹂韩王曰。﹁善。﹂乃警公仲之行。将西购于秦。楚王闻之大恐。召陈轸告之。陈轸曰。﹁秦之欲伐楚久矣。今又得韩之名都一而具甲。秦韩并兵而伐楚。此秦所祷祀而求也。今已得之矣。楚国必伐矣。王听臣为之警四境之内。起师言救韩。命战车满道路。发信臣。多其车。重其币。使信王之救己也。纵韩不能听我。韩必德王也。必不为雁行以来。是秦韩不和也。兵虽至。楚不大病也。为能听我绝和于秦。秦必大怒。以厚怨韩。韩之南交楚。必轻秦。轻秦。其应秦必不敬。是因秦。韩之兵而免楚国之患也。﹂楚王曰。﹁善。﹂乃警四境之内。兴师言救韩。命战车满道路。发信臣。多其车。重其币。谓韩王曰。﹁不谷国虽小。已悉发之矣。愿大国遂肆志于秦。不谷将以楚殉韩。﹂韩王闻之大说。乃止公仲之行。公仲曰。﹁不可。夫以实伐我者秦也。以虚名救我者楚也。王恃楚之虚名。而轻绝强秦之敌。王必为天下大笑。且楚韩非兄弟之国也。又非素约而谋伐秦也。已有伐形。因发兵言救韩。此必陈轸之谋也。且王已使人报于秦矣。今不行。是欺秦也。夫轻欺强秦而信楚之谋臣。恐王必悔之。﹂韩王不听。遂绝于秦。秦因大怒。益甲伐韩。大战。楚救不至韩。十九年。大破我岸门。太子仓质于秦以和。\\
\hspace*{2em}二十一年。与秦共攻楚。败楚将屈丐。斩首八万于丹阳。是岁。宣惠王卒。太子仓立。是为襄王。\\
\hspace*{2em}襄王四年。与秦武王会临晋。其秋。秦使甘茂攻我宜阳。五年。秦拔我宜阳。斩首六万。秦武王卒。六年。秦复与我武遂。九年。秦复取我武遂。十年。太子婴朝秦而归。十一年。秦伐我。取穰。与秦伐楚。败楚将唐眛。\\
\hspace*{2em}十二年。太子婴死。公子咎。公子虮虱争为太子。时虮虱质于楚。苏代谓韩咎曰。﹁虮虱亡在楚。楚王欲内之甚。今楚兵十余万在方城之外。公何不令楚王筑万室之都雍氏之旁。韩必起兵以救之。公必将矣。公因以韩楚之兵奉虮虱而内之。其听公必矣。必以楚韩封公也。﹂韩咎从其计。\\
\hspace*{2em}楚围雍氏。韩求救于秦。秦未为发。使公孙昧入韩。公仲曰。﹁子以秦为且救韩乎。﹂对曰。﹁秦王之言曰﹃请道南郑。蓝田。出兵于楚以待公﹄。殆不合矣。﹂公仲曰。﹁子以为果乎。﹂对曰。﹁秦王必祖张仪之故智。楚威王攻梁也。张仪谓秦王曰。﹃与楚攻魏。魏折而入于楚。韩固其与国也。是秦孤也。不如出兵以到之。魏楚大战。秦取西河之外以归。﹄今其状阳言与韩。其实阴善楚。公待秦而到。必轻与楚战。楚阴得秦之不用也。必易与公相支也。公战而胜楚。遂与公乘楚。施三川而归。公战不胜楚。楚塞三川守之。公不能救也。窃为公患之。司马庚三反于郢。甘茂与昭鱼遇于商于。其言收玺。实类有约也。﹂公仲恐。曰。﹁然则奈何。﹂曰。﹁公必先韩而后秦。先身而后张仪。公不如亟以国合于齐楚。齐楚必委国于公。公之所恶者张仪也。其实犹不无秦也。﹂于是楚解雍氏围。\\
\hspace*{2em}苏代又谓秦太后弟芈戎曰。﹁公叔。伯婴恐秦楚之内虮虱也。公何不为韩求质子于楚。楚王听入质子于韩。则公叔。伯婴知秦楚之不以虮虱为事。必以韩合于秦楚。秦楚挟韩以窘魏。魏氏不敢合于齐。是齐孤也。公又为秦求质子于楚。楚不听。怨结于韩。韩挟齐魏以围楚。楚必重公。公挟秦楚之重以积德于韩。公叔。伯婴必以国待公。﹂于是虮虱竟不得归韩。韩立咎为太子。齐。魏王来。\\
\hspace*{2em}十四年。与齐。魏王共击秦。至函谷而军焉。十六年。秦与我河外及武遂。襄王卒。太子咎立。是为厘王。\\
\hspace*{2em}厘王三年。使公孙喜率周。魏攻秦。秦败我二十四万。虏喜伊阙。五年。秦拔我宛。六年。与秦武遂地二百里。十年。秦败我师于夏山。十二年。与秦昭王会西周而佐秦攻齐。齐败。愍王出亡。十四年。与秦会两周间。二十一年。使暴䳒救魏。为秦所败。䳒走开封。\\
\hspace*{2em}二十三年。赵。魏攻我华阳。韩告急于秦。秦不救。韩相国谓陈筮曰。﹁事急。愿公虽病。为一宿之行。﹂陈筮见穰侯。穰侯曰。﹁事急乎。故使公来。﹂陈筮曰。﹁未急也。﹂穰侯怒曰。﹁是可以为公之主使乎。夫冠盖相望。告敝邑甚急。公来言未急。何也。﹂陈筮曰。﹁彼韩急则将变而佗从。以未急。故复来耳。﹂穰侯曰。﹁公无见王。请今发兵救韩。﹂八日而至。败赵。魏于华阳之下。是岁。厘王卒。子桓惠王立。\\
\hspace*{2em}桓惠王元年。伐燕。九年。秦拔我陉。城汾旁。十年。秦击我于太行。我上党郡守以上党郡降赵。十四年。秦拔赵上党。杀马服子卒四十余万于长平。十七年。秦拔我阳城。负黍。二十二年。秦昭王卒。二十四年。秦拔我城皋。荥阳。二十六年。秦悉拔我上党。二十九年。秦拔我十三城。\\
\hspace*{2em}三十四年。桓惠王卒。子王安立。\\
\hspace*{2em}王安五年。秦攻韩。韩急。使韩非使秦。秦留非。因杀之。\\
\hspace*{2em}九年。秦虏王安。尽入其地。为颍川郡。韩遂亡。\\
\hspace*{2em}太史公曰。韩厥之感晋景公。绍赵孤之子武。以成程婴。公孙杵臼之义。此天下之阴德也。韩氏之功。于晋未睹其大者也。然与赵。魏终为诸侯十余世。宜乎哉。\\

% ===== 索隐述赞（司马贞《史记索隐》）· 统一样式 =====
% 题头「索隐述赞曰」五字并入反白框；正文每四句一列，定宽盒 \hfil 居列中；
% \fontsize{24}{41} 列距对齐全书界行（121+41k）；\clearpage 使述赞独占一页。
\clearpage
\shuzanhead{索隐述赞曰}\linebreak
{\shuzanbody\fontsize{24}{41}\selectfont
\hbox to \hsize{\hfil 韩氏之先。实宗周武。事微国小。春秋无语。\hfil}\linebreak
\hbox to \hsize{\hfil 後裔事晋。韩原是处。赵孤克立。智伯可取。\hfil}\linebreak
\hbox to \hsize{\hfil 既徙平阳。又侵负黍。景赵俱侯。惠又僭主。\hfil}\linebreak
\hbox to \hsize{\hfil 秦败脩鱼。魏会区鼠。韩非虽使。不禁狼虎。\hfil}%
\par}
\end{document}
